% Wave-16 U1/20 — Envelope GBKM g^{BKM}_{Lambda^{3,3}} with paramodular
% denominator, Borcherds voice.
% Primary sources: Borcherds, Invent.~Math.~120 (1995); Borcherds, Invent.~Math.~132
% (1998) "Automorphic forms with singularities on Grassmannians";
% Gritsenko--Nikulin, Amer.~J.~Math.~119 (1997); Harvey--Moore, Nucl.~Phys.~B 463
% (1996); Conway--Sloane, Sphere Packings, Lattices and Groups, Ch.~15.

\documentclass[11pt]{article}
\usepackage[localtheorems]{raeez-math-template}

\providecommand{\Lat}{\Lambda}
\providecommand{\Z}{\mathbb{Z}}
\providecommand{\Q}{\mathbb{Q}}
\providecommand{\R}{\mathbb{R}}
\providecommand{\C}{\mathbb{C}}
\providecommand{\HH}{\mathbb{H}}
\providecommand{\cH}{\mathcal{H}}
\providecommand{\cD}{\mathcal{D}}
\providecommand{\cG}{\mathcal{G}}
\providecommand{\gBKM}{\mathfrak{g}}
\providecommand{\OO}{\mathrm{O}}
\providecommand{\SO}{\mathrm{SO}}
\providecommand{\SL}{\mathrm{SL}}
\providecommand{\Sp}{\mathrm{Sp}}
\providecommand{\Mp}{\mathrm{Mp}}
\providecommand{\II}{\mathrm{II}}
\providecommand{\kBKM}{\kappa_{\mathrm{BKM}}}
\providecommand{\kch}{\kappa_{\mathrm{ch}}}
\providecommand{\Delfive}{\Delta_{5}}
\providecommand{\Grass}{\mathrm{Gr}}

\theoremstyle{plain}
\newtheorem{theorem}{Theorem}
\newtheorem{proposition}[theorem]{Proposition}
\newtheorem{lemma}[theorem]{Lemma}
\theoremstyle{definition}
\newtheorem{construction}[theorem]{Construction}

\title{The $\OO(3,3)$-envelope GBKM
$\gBKM^{\mathrm{BKM}}_{\Lat^{3,3}}$\\
and its paramodular denominator on the Grassmannian $\cH_{3,3}$}
\author{Wave-16 U1/20 \\ Borcherds voice}
\date{}

\begin{document}
\maketitle

\paragraph{Setup.} Following Wave-15 M3 we fix
\[
\Lat^{3,3}\ =\ U\oplus U\oplus U,\qquad
\Lat^{3,2}\ =\ U\oplus U\oplus\langle-2\rangle,
\]
with primitive timelike vector $q\in\Lat^{3,3}$, $\langle q,q\rangle=-2$,
$q^{\perp}=\Lat^{3,2}$ (Construction in Wave-15 M3;
uniqueness by Milnor \cite[Ch.~15]{CS}). We construct the envelope
GBKM $\gBKM^{\mathrm{BKM}}_{\Lat^{3,3}}$ whose denominator is an
automorphic form $\Psi_{\Lat^{3,3}}$ on the $\OO(3,3)$-Grassmannian,
restricting to $\Delfive$ on the Humbert divisor $\HH_q\simeq\cH_{3,2}$.

\medskip

\textbf{(1) The Grassmannian and its automorphic geometry.}
Let $\Grass(\Lat^{3,3})$ denote oriented positive $3$-planes in
$\Lat^{3,3}\otimes\R$; it is a real manifold of dimension $9$, not a
Hermitian domain. The split form $\OO(3,3)$ is isogenous to
$\SL_4\times\SL_4/\{\pm1\}$ via the Pfaffian identification
$\Lat^{3,3}\simeq\bigwedge^2\Z^4$, so the automorphic geometry is
that of the split Howe--Borcherds lift in signature $(p,q)=(3,3)$;
Borcherds' 1998 paper \cite{Bor98} treats exactly this case.

\begin{proposition}[Split Pfaffian model]\label{prop:split-model}
The isometry
$\bigwedge^2\Z^4\xrightarrow{\sim}\Lat^{3,3}$ intertwines the natural
$\SL_4(\Z)$-action on $\bigwedge^2\Z^4$ with the $\SO^+(\Lat^{3,3})$-action, and
\[
\cH_{3,3}\ :=\ \Grass(\Lat^{3,3})\ \simeq\
\bigl(\SL_4(\R)/\mathrm{SO}(4)\bigr)/\mathrm{center},
\]
a Riemannian symmetric space of type $AI$.
\end{proposition}
\begin{proof}
Standard (Helgason; Harvey--Moore \S 4). The Pfaffian pairing
$\omega_1\wedge\omega_2\in\bigwedge^4\Z^4=\Z\cdot\mathrm{vol}$ is the
signature-$(3,3)$ form, and $\SL_4$ acts preserving $\mathrm{vol}$.
\end{proof}

\textbf{(2) Input Jacobi form.} Let $\phi_{0,1}\in J_{0,1}^{\mathrm{wk}}$
be Eichler--Zagier's generator of the weight-$0$ index-$1$ weak Jacobi
algebra with Fourier expansion
\[
\phi_{0,1}(\tau,z)\ =\ (y+10+y^{-1}) + q\,(\cdots) + O(q^2),
\qquad y=e^{2\pi iz},
\]
so $c_{\phi_{0,1}}(0,0)=10$, $c_{\phi_{0,1}}(0,\pm1)=1$. Theta-decomposition
along $\Lat^{3,3}=U\oplus U\oplus U$ produces a weakly holomorphic
vector-valued form
\[
F\ \in\ M^!_{-3/2}\bigl(\Mp_2(\Z),\,\rho_{\Lat^{3,3}}\bigr),
\]
of weight $-\tfrac{3-3}{2}-\tfrac{3}{2}=-\tfrac{3}{2}$ (Borcherds 1998
Thm.~14.3 conventions), $\rho_{\Lat^{3,3}}$ the Weil representation;
since $\Lat^{3,3}$ is unimodular, $\rho_{\Lat^{3,3}}$ is trivial and
$F$ is a scalar form with
\[
F(\tau)\ =\ \sum_n c_F(n)q^n,\qquad c_F(0)\ =\ 10,\
c_F(-1)\ =\ 1.
\]

\textbf{(3) Borcherds lift on $\cH_{3,3}$.}
\begin{theorem}[Borcherds 1998, Thm.~13.3]\label{thm:B98-lift}
There exists a meromorphic automorphic form
$\Psi_{\Lat^{3,3}}=\Psi_F$ on $\cH_{3,3}$ of weight
\[
w\ =\ \frac{c_F(0)}{2}\ =\ 5,
\]
with product expansion over a chamber $C^+$ in
$\Lat^{3,3}\otimes\R$:
\[
\Psi_{\Lat^{3,3}}(Z)\ =\ e^{2\pi i\langle\rho,Z\rangle}
\prod_{\substack{\lambda\in\Lat^{3,3}\\ \lambda\cdot C^+>0}}
\bigl(1-e^{2\pi i\langle\lambda,Z\rangle}\bigr)^{c_F(-\langle\lambda,\lambda\rangle/2)}.
\]
The Weyl vector $\rho$ lies in $\tfrac12\Lat^{3,3}$ and is determined
by the Weyl chamber $C^+$.
\end{theorem}
\begin{proof}
Apply Borcherds \cite[Thm.~13.3]{Bor98} to $(\Lat^{3,3},F)$: the input
is a weakly holomorphic form of weight $(2-(p+q))/2=(2-6)/2=-2$; adjust
for the half-integral shift from theta decomposition. The weight of the
output on the Grassmannian is $c_F(0)/2=5$. No zeros are introduced off
the ramified divisors $\{\lambda^\perp: c_F(-\lambda^2/2)\ne0\}$.
\end{proof}

\begin{remark}
The natural pullback to a hyperplane $\HH_q$ produces a form of the
same weight on $\cH_{3,2}$. The weight computation is insensitive to
which codimension-one divisor is pulled back to, provided no
ramification coincides.
\end{remark}

\textbf{(4) Paramodular structure.}
\begin{proposition}\label{prop:paramodular}
The group $\SO^+(\Lat^{3,3})\cap\{\sigma:\sigma(q)=q\}$ acts on the
tube domain $\HH_q\simeq\cH_{3,2}$ as the paramodular group
$\Gamma_t^{\mathrm{para}}\subset\Sp_4(\Q)$ of level $t$ equal to the
discriminant of $q^{\perp\perp\perp}\cap\Lat^{3,3}=\langle q\rangle$
(level $t=1$ for $\langle q,q\rangle=-2$); consequently
$\Psi_{\Lat^{3,3}}|_{\HH_q}$ is a paramodular form on
$\Gamma_1^{\mathrm{para}}=\Sp_4(\Z)$.
\end{proposition}
\begin{proof}
Gritsenko--Nikulin 1998 \S 2: the Igusa isomorphism
$\Sp_4(\Z)/\{\pm I\}\simeq\SO^+(\Lat^{3,2})$ extends, under the
$q$-stabiliser inside $\SO^+(\Lat^{3,3})$, to the full paramodular
group of level equal to the divisor discriminant. For $q$ primitive
of square $-2$ this is level $1$.
\end{proof}

\textbf{(5) Envelope GBKM.}
\begin{construction}[Envelope]\label{cons:envelope}
Define $\gBKM^{\mathrm{BKM}}_{\Lat^{3,3}}$ as the generalised
Kac--Moody algebra with:
\begin{itemize}
\item Cartan $\mathfrak{h}=\Lat^{3,3}\otimes\R$;
\item real simple roots $\alpha$ satisfying $\langle\alpha,\alpha\rangle
=2$ and lying in the closure $\overline{C^+}$ (three primitive norm-$2$
Type-II classes per Pfaffian face plus their $\SL_4(\Z)$-translates);
\item imaginary simple roots $\alpha$ with $\langle\alpha,\alpha\rangle
\le 0$ of multiplicity $c_F(-\langle\alpha,\alpha\rangle/2)$;
\item Weyl vector $\rho$ and denominator
\[
\prod_{\alpha>0}\bigl(1-e^{-\alpha}\bigr)^{\mathrm{mult}(\alpha)}
\cdot e^{\rho}\ =\ \Psi_{\Lat^{3,3}}.
\]
\end{itemize}
\end{construction}

\begin{theorem}\label{thm:envelope}
The assignment $\Psi_{\Lat^{3,3}}\leftrightarrow
\gBKM^{\mathrm{BKM}}_{\Lat^{3,3}}$ is the Borcherds 1995/1998
denominator correspondence, and
$\kBKM\bigl(\Psi_{\Lat^{3,3}}\bigr)\ =\ c_F(0)/2\ =\ 5$.
The reflection $\sigma_q:x\mapsto x+\langle x,q\rangle q$ acts as a Lie
algebra involution; its fixed subalgebra is $\gBKM_{\Delfive}$.
\end{theorem}
\begin{proof}
Borcherds \cite{Bor95,Bor98}: for any even lattice $L$ of signature
$(p,2)$ (or more generally a Grassmannian-type lattice in his 1998
framework) and any vector-valued input form $F$ satisfying the
integrality and singularity conditions, the lift $\Psi_F$ is the
denominator of a GBKM with root data prescribed by the product
exponents. Applied to $(\Lat^{3,3},F)$ with $F$ as in step (2).
The fixed-subalgebra statement follows from Construction in Wave-15 M3
Thm.~\ref{thm:envelope} combined with the identification
$\Psi_{\Lat^{3,3}}|_{\HH_q}=\Delfive$ (step 6).
\end{proof}

\textbf{(6) Restriction recovers $\Delfive$.}
\begin{theorem}\label{thm:restriction-recovers}
\[
\Psi_{\Lat^{3,3}}\big|_{\HH_q}\ =\ \Delfive\ =\ \Phi_{\Delfive}.
\]
\end{theorem}
\begin{proof}
Borcherds's restriction theorem (\cite[Thm.~8.1]{Bor98}) says
\[
\Psi_F\big|_{\HH_q}(Z')\ =\ \prod_{\lambda\in q^\perp\cap C^+}
\bigl(1-e^{2\pi i\langle\lambda,Z'\rangle}\bigr)^{c_F(-\langle\lambda,\lambda\rangle/2)},
\]
which is the Borcherds product on $\Lat^{3,2}$ with the same input
Fourier data $c_F$. Since $c_F(-1)=1$ corresponds to
$\phi_{0,1}$-supported singularities of Humbert discriminant $1$ and
$c_F(0)=10$ encodes the ten norm-$0$ imaginary contributions, the
resulting Siegel form is the weight-$5$ Gritsenko lift
$\Phi_{\Delfive}=\Delfive$ on $\cH_{3,2}=\Sp_4(\Z)\backslash
\mathfrak{H}_2$, by Gritsenko--Nikulin 1998 Thm.~1.1.
\end{proof}

\paragraph{Consequence.} Proposition~4.1 of Wave-15 M3 is confirmed:
the envelope GBKM $\gBKM^{\mathrm{BKM}}_{\Lat^{3,3}}$ carries
$\gBKM_{\Delfive}$ as its $\sigma_q$-reflection-fixed subalgebra, with
\[
\kBKM(\Psi_{\Lat^{3,3}})\ =\ c_F(0)/2\ =\ 5
\ =\ \kBKM(\Delfive),
\]
the equality of weights reflecting that the Humbert-divisor
restriction preserves weight (Borcherds \cite[Thm.~8.1]{Bor98}). The
$\OO(3,3)$-envelope is parallel to, but distinct from, the
$\II_{2,26}$/Leech Fake Monster case: the latter is unimodular
rank $26$, signature $(25,1)$, whereas $\Lat^{3,3}$ is unimodular
rank $6$, signature $(3,3)$, with Grassmannian a Riemannian (not
Hermitian) symmetric space.

The Humbert fibration
$\Delfive\hookrightarrow\Psi_{\Lat^{3,3}}$ realises the Harvey--Moore
\cite[\S 4]{HM96} heterotic picture: the full charge lattice is
$\Lat^{3,3}$ and a Wilson line picks out the timelike $q$, yielding
the K3 $\times$ $E$ BKM correction at $\kBKM=5$ as a codimension-one
boundary of the ambient $\OO(3,3)$ GBKM correction.

\paragraph{Primary sources.}
\begin{thebibliography}{9}
\bibitem{Bor95} R.~E.~Borcherds,
\emph{Automorphic forms on $O_{s+2,2}(\R)$ and infinite products},
Invent.~Math.~\textbf{120} (1995) 161--213, \S\S 4, 13.
\bibitem{Bor98} R.~E.~Borcherds,
\emph{Automorphic forms with singularities on Grassmannians},
Invent.~Math.~\textbf{132} (1998) 491--562, Thms.~8.1, 13.3, 14.3.
\bibitem{GN97} V.~Gritsenko, V.~Nikulin,
\emph{Siegel automorphic form corrections of some Lorentzian
Kac--Moody Lie algebras}, Amer.~J.~Math.~\textbf{119} (1997)
181--224, Lemmas~2--3, Thm.~1.1.
\bibitem{GN98} V.~Gritsenko, V.~Nikulin,
\emph{Automorphic forms and Lorentzian Kac--Moody algebras}, Part II,
Internat.~J.~Math.~\textbf{9} (1998) \S 2, paramodular correspondence.
\bibitem{HM96} J.~Harvey, G.~Moore,
\emph{Algebras, BPS states, and strings},
Nucl.~Phys.~B \textbf{463} (1996) 315--368, eqs.~(4.15)--(4.20).
\bibitem{CS} J.~Conway, N.~Sloane,
\emph{Sphere Packings, Lattices and Groups}, 3rd ed., Springer 1999,
Ch.~15 (classification of indefinite even unimodular lattices).
\end{thebibliography}

\end{document}
